% =====================================================================
\section*{Question 1 -- River Forecasting Proposal}
\addcontentsline{toc}{section}{Question 1}

% ---------------------------------------------------------------------
\subsection*{1. Architecture diagram}

All sources are placed on one space--time grid (places $\times$ ten-minute steps). A causal temporal convolution summarises what has happened at each place. Then, four rounds of upstream attention pass information down the river, with a learned travel delay. A small head predicts each station's change in level over the next hour.

\begin{figure}[H]\centering
\includegraphics[width=\linewidth,height=0.59\textheight,keepaspectratio]{../Question1/figures/river_architecture.pdf}
\caption{Proposed forecasting system.}
\label{fig:q1arch}
\end{figure}

% ---------------------------------------------------------------------
\subsection*{2. Assumptions}

Before stating assumptions, we list the facts given in the brief.
\begin{itemize}[label={},leftmargin=0.6em]
  \item \textbf{F1} levels at S1--S4 and T1 every 10 minutes;
  \item \textbf{F2} rainfall at R1, R2 with timestamps and gauge locations;
  \item \textbf{F3} recorded releases, and future release schedules with announcement and effective times;
  \item \textbf{F4} station locations, connectivity and river distances;
  \item \textbf{F5} a monitored tributary (T1) joins between S1 and S2, an unmonitored one between S2 and S3;
  \item \textbf{F6} stations are unevenly spaced and travel times change with conditions;
  \item \textbf{F7} future rainfall, future levels, exact travel times and the unmonitored inflow are not available.
\end{itemize}

Based on what the brief leaves unspecified, we make the following assumptions. Source markers [1]--[6] refer to the list below the table.

{\small
\begin{xltabular}{\linewidth}{@{}p{0.55cm}LLL@{}}
\caption{Assumptions, why each is reasonable, and its effect on the design.}\label{tab:q1assump}\\
\toprule
 & Assumption & Why reasonable & Effect on design\\
\midrule
\endfirsthead
\toprule
 & Assumption & Why reasonable & Effect on design\\
\midrule
\endhead
\bottomrule
\endlastfoot
\multicolumn{4}{@{}l}{\textbf{Data assumptions} \ (record, timing and quality the brief leaves open)}\\
A1 & Multi-year record with several high-water events. & Deep hourly rainfall--runoff models may need $\sim$15 years to beat simpler models [4]. & Chronological split by year; small model (A16); always compared with simple baselines (Section 4).\\
A2 & One clock (UTC); level readings are instantaneous values. & Standard telemetry practice. & Snap to nearest 10-min mark ($\pm$2 min), else missing.\\
A3 & Each rain reading is the total since the previous one. & Tipping-bucket gauges tip at a fixed depth ($\sim$0.25 mm) [2]. & Sum into 10-min bins.\\
A4 & Readings are recorded every 10 min (F1) but sent in batches, at least hourly; the newest can be up to $\sim$1 h old at issuance. & Gauges commonly record every 15 min and transmit about hourly, more often in floods [1]. & Issue on time; readings not yet received are missing; training simulates delays up to 1 h (Section 5).\\
A5 & Data contain gaps, spikes and flat-lines, but no long-term drift. & Common sensor faults. & Quality screen; \emph{observed?} flag per cell.\\
A6 & Release schedules are archived with announcement and effective times. & The authority records announcements (F3). & Train on plans \emph{as announced} (no leakage).\\
A7 & Dam-to-S4 travel $\le\sim$4 h at low flow; rain response mostly within $\sim$5 h. & Short mountain river (handout figure); checked by the lag of best upstream--downstream correlation. & History window $H=30$ steps (5 h).\\
A8 & Each station has alert levels $\theta_s$ (or high percentiles of the record). & Authorities define action/minor/moderate/major flood stages [3]. & Larger loss weight above $\theta_s$.\\
\midrule
\multicolumn{4}{@{}l}{\textbf{River-dynamics assumptions} \ (how the model treats the river)}\\
A9 & Water moves only downstream; no backwater effect within an hour. & Steep river with free-flowing junctions. & Attention restricted to upstream places (generic spatial blocks look both ways).\\
A10 & Travel time is positive and varies with flow: usually shorter as the river rises in its channel, but can lengthen once water spreads onto the floodplain. & Wave speed $\approx\tfrac53\times$ flow velocity in a channel [5], but may fall above bankfull where there is floodplain storage [6]. & Delay learned by state-dependent attention, not fixed.\\
A11 & R1 drains to the S1--S2 reach; R2 to the unmonitored tributary (to S3). & Gauge positions in the handout figure. & R1 added to row S2, R2 to row S3.\\
A12 & Unmonitored inflow is driven mainly by rain on its catchment, plus slow baseflow. & Typical of small tributaries. & R2 as proxy; remaining error widens uncertainty at S3/S4.\\
A13 & Releases are the only controlled inflow; withdrawals are negligible. & None mentioned. & No extra input rows.\\
A14 & Without a new announcement, the current release continues. & Operators change releases by announcement. & Future Dam cells hold the current rate.\\
A15 & Gauge datums and channel shape are stable over the record. & No changes mentioned. & One model for the whole record; retrain if changed.\\
A16 & The same learned update rule applies at every station; reaches differ only through place embedding and distance. & Few stations and limited flood data; sharing pools examples. & Shared TCN and attention weights.\\
\end{xltabular}}

{\footnotesize\noindent\textbf{Sources.}
[1] USGS, \href{https://usgs.gov/faqs/how-often-are-real-time-streamflow-data-updated}{How often are real-time streamflow data updated?};
[2] AMS Glossary, \href{https://glossary.ametsoc.org/wiki/Tipping-bucket_rain_gauge}{Tipping-bucket rain gauge};
[3] NWS, \href{https://weather.gov/lot/hydrology_definitions}{Hydrology definitions};
[4] \href{https://pmc.ncbi.nlm.nih.gov/articles/PMC12055974}{Evaluation of LSTM vs.\ conceptual models for hourly rainfall--runoff simulations with varied training period lengths} (2025);
[5] Colorado State Univ., \href{https://www.engr.colostate.edu/~pierre/ce_old/classes/ce717/Chapter\%2011\%20Unsteady\%20Flow\%20in\%20Open\%20Channels.pdf}{CE717 notes, Unsteady flow in open channels};
[6] \href{https://meetingorganizer.copernicus.org/EGU2017/EGU2017-9406.pdf}{EGU General Assembly 2017, abstract EGU2017-9406} (wave celerity vs.\ discharge).\par}

% ---------------------------------------------------------------------
\subsection*{3. Forecasting task and data specification}

\textbf{Issuance and horizon.} A forecast is issued every 10 minutes at time $t_0$ (:00, :10, \dots). It predicts the level at S1--S4 at lead times $t_0+10k$ min, $k=1,\dots,6$: a $4\times6$ output per issuance.

\textbf{History window.} The model sees the last $H=30$ steps (5 h, $t_0-290$ min to $t_0$), which covers the assumed travel and rainfall-response times (A7). The delay within the window is learned.

\textbf{Grid.} All inputs share one grid: $T=H+6=36$ columns ($\tau=-29,\dots,0$ past; $\tau=1,\dots,6$ future) and $P=6$ rows (Dam, S1, T1, S2, S3, S4). Each cell holds $F=4$ features $[\ell, m, r, q]$: level, observed flag, rain and release. $m=1$ only if every measured value in the cell has arrived. In the Dam row, past columns hold actual releases and future columns the announced plan; column position separates the two.

\begin{table}[H]\centering\small
\caption{Inputs and outputs for one batch. $B$ = number of issuances (forecast times) in the batch.}
\label{tab:q1data}
\begin{tabularx}{\linewidth}{@{}l>{\raggedright\arraybackslash}p{2.3cm}LL@{}}
\toprule
Item & Shape & Axes & Time coverage\\
\midrule
Grid $\mathbf X$ & $B\times36\times6\times4$ & issuance, column $\tau$, place, feature & $t_0-290$ min \dots\ $t_0+60$ min\\
\quad level $\ell$ & (in $\mathbf X$) & rows S1--S4, T1 & $\tau\le0$ as observed; $\tau>0$ empty\\
\quad observed flag $m$ & (in $\mathbf X$) & all rows & 1 if every measured value has arrived; 0 otherwise and for $\tau>0$\\
\quad rain $r$ & (in $\mathbf X$) & R1 in row S2, R2 in row S3 & $\tau\le0$; $\tau>0$ empty\\
\quad release $q$ & (in $\mathbf X$) & Dam row & $\tau\le0$ actual; $\tau>0$ plan announced by $t_0$\\
Calendar $\mathbf C$ & $B\times36\times4$ & issuance, column, $\sin/\cos$ of time of day and day of year & same as $\mathbf X$\\
River map $\mathcal E$ & 5 edges, each with distance $d_{ji}$ & Dam$\to$S1, S1$\to$S2, T1$\to$S2, S2$\to$S3, S3$\to$S4 & fixed\\
Output $\hat{\mathbf Y}$ & $B\times4\times6$ & issuance, station S1--S4, lead time $k$ & $t_0+10$ \dots\ $t_0+60$ min\\
Target $\mathbf Y$, mask $\mathbf M$ & $B\times4\times6$ & as $\hat{\mathbf Y}$ (training only) & same\\
\bottomrule
\end{tabularx}
\end{table}

\textbf{Alignment across sampling times.} All sources use one 10-minute clock. A level reading enters column $\tau$ if its timestamp lies within $\pm$2 min of $t_0+10\tau$; otherwise the cell is missing. Rain totals (A3) are summed over $(t_0+10(\tau-1),\,t_0+10\tau]$ by observation time. Releases are averaged per interval and placed by effective time.

\textbf{Alignment across locations.} Values are not interpolated in space. Each source maps to one place row (rain via A11). River distances enter as edge features in Pass Downstream, so unequal spacing is handled by distances and learned delays, not by treating rows as evenly spaced neighbours.

\textbf{Information available at issuance.} A cell is filled only if its value is known at $t_0$: a measurement must have arrived, and a plan must have been announced, by $t_0$. Future levels and rainfall are never available; their cells stay empty with $m=0$. The only filled future cells are the Dam row's announced releases.

% ---------------------------------------------------------------------
\subsection*{4. Model computations and design justification}

\textbf{Notation.} Places $p\in\{\text{Dam},\text{S1},\text{T1},\text{S2},\text{S3},\text{S4}\}$; columns $\tau\in\{-29,\dots,6\}$; stations $s\in\{\text{S1},\dots,\text{S4}\}$; lead times $k=1,\dots,6$; hidden width $d=64$. $\mathbf x_{p,\tau}=[\ell,m,r,q]\in\mathbb R^4$ is one grid cell, $\mathcal U(i)$ the set of places directly upstream of $i$ (from $\mathcal E$), and $d_{ji}$ the river distance from $j$ to $i$. Levels are standardised per station using training statistics.

\textbf{Input projection.}
\begin{equation}
\mathbf z^{(0)}_{p,\tau} = W_{\text{in}}\,\mathbf x_{p,\tau} + \mathbf e_p + \mathbf c_\tau, \qquad \mathbf z\in\mathbb R^{d},
\end{equation}
where $\mathbf e_p$ is a learned place embedding and $\mathbf c_\tau$ a time embedding (calendar features plus column position, so future columns are distinguishable).

\textbf{Look Back (Causal Dilated TCN), applied to each place with shared weights.} For layers $l=1,\dots,5$ with kernel size 3 and dilation $\delta_l=2^{l-1}$:
\begin{equation}
\mathbf z^{(l)}_{p,\tau} = \mathbf z^{(l-1)}_{p,\tau} + \mathrm{ReLU}\Big(\textstyle\sum_{a=0}^{2} A^{(l)}_a\,\mathbf z^{(l-1)}_{p,\,\tau-a\delta_l} + \mathbf b^{(l)}\Big),
\end{equation}
with zero padding before $\tau=-29$. The convolution is causal (columns $\le\tau$ only). Its receptive field is $1+2(1+2+4+8+16)=63\ge36$ columns, so each column sees the full window. Set $\mathbf h^{0}_{p,\tau}=\mathbf z^{(5)}_{p,\tau}$.

\textbf{Pass Downstream (Upstream Attention $\times4$), $R=4$ rounds with separate weights per round.} For round $r$ and each place $i$ with $\mathcal U(i)\neq\emptyset$:
\begin{align}
\mathbf q_{i,\tau} &= W_Q^{(r)}\mathbf h^{r-1}_{i,\tau}, \qquad
\mathbf k_{j,\tau'} = W_K^{(r)}\mathbf h^{r-1}_{j,\tau'} + \mathbf u^{(r)}_{\tau-\tau'} + \mathbf w^{(r)} d_{ji}, \qquad
\mathbf v_{j,\tau'} = W_V^{(r)}\mathbf h^{r-1}_{j,\tau'},\\
\alpha_{i\tau,\,j\tau'} &= \operatorname*{softmax}_{\,j\in\mathcal U(i),\ \tau'\le\tau}\!\Big(\mathbf q_{i,\tau}^{\top}\mathbf k_{j,\tau'}/\sqrt d\Big), \qquad
\mathbf m_{i,\tau} = \textstyle\sum_{j,\tau'}\alpha_{i\tau,\,j\tau'}\,\mathbf v_{j,\tau'},\\
\tilde{\mathbf h}_{i,\tau} &= \mathrm{LN}\big(\mathbf h^{r-1}_{i,\tau} + W_O^{(r)}\mathbf m_{i,\tau}\big), \qquad
\mathbf h^{r}_{i,\tau} = \mathrm{LN}\big(\tilde{\mathbf h}_{i,\tau} + \mathrm{FFN}^{(r)}(\tilde{\mathbf h}_{i,\tau})\big).
\end{align}
$\mathbf u_{\tau-\tau'}$ is a learned lag embedding and $\mathbf w d_{ji}$ a distance term. Places without upstream neighbours (Dam, T1) keep $\mathbf h^{r}=\mathbf h^{r-1}$. Four heads are used; one is written for clarity.

\textbf{Predict Ahead (MLP).} Each station reads its own representation at future column $k$:
\begin{equation}
\Delta_{s,k} = g\big(\mathbf h^{4}_{s,k}\big), \qquad \hat y_{s,k} = \tilde\ell_{s,0} + \sigma_s\,\Delta_{s,k},
\end{equation}
where $g$ is a two-layer MLP shared by all stations (the place embedding separates them), $\tilde\ell_{s,0}$ is the latest valid level at $s$, and $\sigma_s$ the training standard deviation of one-hour level changes at $s$.

\textbf{Training objective.} Weighted, masked mean squared error, as in the tutorial:
\begin{equation}
\mathcal L = \frac{\sum_{b,s,k} M_{b,s,k}\,w_{b,s,k}\,\big(\hat y_{b,s,k}-y_{b,s,k}\big)^2/\sigma_s^2}{\sum_{b,s,k} M_{b,s,k}\,w_{b,s,k}}, \qquad
w_{b,s,k} = 1 + \lambda\,\mathbb 1\big[y_{b,s,k}\ge\theta_s\big],
\end{equation}
with $M=0$ for missing targets, $\lambda=4$ (tuned on validation) and alert levels $\theta_s$ (A8). Training uses Adam, early stopping on validation error and a chronological split (A1). 10--20\% of cells, and occasionally a whole row, are masked at random so the model learns to forecast with missing data. MSE targets the conditional mean, so peaks are smoothed when information is missing; the high-water weight counters this (Table~\ref{tab:q1limits}). \emph{Optional extension:} $g$ outputs the 10/50/90\% quantiles, trained with pinball loss, to give a forecast range.

\textbf{Why this design.}
\begin{itemize}
  \item \emph{Upstream paths are required.} A station's own history cannot show water already in transit (companion notebook, Part 1). Pass Downstream gives each station a path to every upstream source, with neighbours defined by the river map (F4), not by index order. Four rounds cover the longest path, Dam$\to$S1$\to$S2$\to$S3$\to$S4.
  \item \emph{Delay must be learned, not fixed.} Travel time changes with flow (F6, A10). Attention weights are recomputed for each input, so the same model can look 50 min back on a calm day, 25 min back as the river rises, and further back if an overbank flood slows the wave. A fixed kernel would impose one delay.
  \item \emph{TCN for local history.} A causal TCN handles the look-back step; an LSTM/GRU would be an equally reasonable choice. The TCN was chosen for its exact receptive field (63 $\ge$ 36 steps), parallel computation, and consistency with the causal temporal blocks in the companion notebook. It also extracts local shape (rising limb, peak), which single-value attention tokens cannot represent.
  \item \emph{Announced releases on the same grid.} Covariates help only when they cover the forecast horizon (PA1). The plan in the Dam row's future cells reaches S1 and onward through the same blocks.
  \item \emph{Predicting the change.} Levels persist over one hour. Predicting $\Delta$ makes persistence the default, so the network learns only rises and falls.
\end{itemize}

\textbf{How each source reaches each forecast.}
\begin{itemize}
  \item \emph{Reservoir releases} (past and planned) enter at the Dam row and reach S1, S2, S3 and S4 in rounds 1--4.
  \item \emph{Upstream conditions}: each station attends to its upstream neighbours, whose states already carry information from further upstream.
  \item \emph{Monitored tributary}: T1 is a direct upstream neighbour of S2, and through S2 of S3 and S4.
  \item \emph{Rainfall}: R1 enters S2's row (local inflow); R2 enters S3's row as the proxy for the unmonitored tributary (A11, A12). The TCN learns the rain-to-river delay.
\end{itemize}

\textbf{Physical assumptions vs.\ learned mappings.}
\begin{table}[H]\centering\small
\begin{tabularx}{\linewidth}{@{}LL@{}}
\toprule
Imposed by design (physical assumptions) & Learned from data\\
\midrule
Flow direction and connectivity: information moves only downstream (A9) & Travel time and how it changes with flow (attention over lags)\\
Causality: no column uses later columns; nothing unknown at $t_0$ & Attenuation and spreading of flood waves between stations\\
Where each source enters the river (Dam$\to$S1, T1$\to$S2, rain rows; A11) & Rainfall-to-level response and its delay\\
Positive, bounded delay: at most 5 h (A7) & Unmonitored inflow, inferred from R2 and the S2--S3 difference\\
Continuity: forecast = current level + change & How reliable announced releases are (plan in future Dam cells, actuals in past cells)\\
\bottomrule
\end{tabularx}
\end{table}
No separate physical simulator is used; physics enters through the structure above. Uncertainty is handled by the optional quantile output and the missing-data treatment in Sections 5 and 6.

\textbf{Baselines for evaluation.} Persistence ($\hat y_{s,k}=\tilde\ell_{s,0}$) and an Autoformer with all places mixed into one token per step (as in PA1). Ablations: no Pass Downstream (each station alone), no future plan, and fixed instead of attended delays.

% ---------------------------------------------------------------------
\subsection*{5. Operation and exceptional conditions}

\textbf{(i) Missing observations.} A missing level or rain value is set to 0 with $m=0$, which separates ``missing'' from ``zero''. Random masking in training (Section 4) teaches the model to rely on upstream and neighbouring rows. If a station's current level is missing, the forecast is anchored on its latest valid value $\tilde\ell_{s,0}$ and marked lower-confidence. Missing targets are excluded from the loss ($M=0$). If a station has been missing for more than two hours, forecasts are still issued, flagged to operators, with a wider range.

\textbf{(ii) Delayed observations.} The forecast is issued on time at $t_0$ with whatever has arrived; values not yet received are missing. A reading is placed by its \emph{observation} timestamp but used only if its \emph{arrival} time is $\le t_0$. A late reading fills its cell at the next issuance; since one forecast takes milliseconds, the current forecast can also be re-issued immediately. Training simulates delays of up to 1 h (A4), matching the gaps seen in operation.

\textbf{(iii) Announced future releases.} The latest schedule announced by $t_0$ fills the Dam row's future columns ($\tau=1,\dots,6$) by effective time. A schedule announced after $t_0$ is not used, even if it takes effect earlier. Without an announcement, the current rate continues (A14). Training uses plans as announced at each past $t_0$ (A6).

\textbf{(iv) Scheduled vs.\ actual releases.} Once a release is recorded, its past cell holds the actual value, which overrides the plan. Training pairs announced plans with the levels that followed, so the model learns how far plans can be trusted. If the actual release departs from the plan beyond a set threshold over the last 30 min, operators are alerted and the range widens. Because a forecast is cheap, operators can also run ``what-if'' forecasts for alternative plans.

\textbf{Feasibility every 10 minutes.} One forecast processes $6\times36=216$ cells. The TCN costs about $216\times5\times3\times64^2\approx1.3\times10^{7}$ multiply--adds; the four attention rounds about $4\times36\times72\times64\times2\times4\approx5\times10^{6}$; the model has about $2\times10^{5}$ parameters. This runs in well under 10 ms on a standard CPU, and alignment needs a few database queries, so a ten-minute cycle leaves a large margin. Training runs offline and is repeated periodically (e.g.\ monthly or after a major flood).

% ---------------------------------------------------------------------
\subsection*{6. Uncertainty, limitations, and additional information}

\begin{table}[H]\centering\small
\caption{Uncertainties and limitations, their effect on forecasts, and our response.}
\label{tab:q1limits}
\begin{tabularx}{\linewidth}{@{}>{\raggedright\arraybackslash}p{2.9cm}LL@{}}
\toprule
Issue & Effect on forecasts & Response / unresolved\\
\midrule
Unmonitored tributary & S3 and S4 under-forecast in storms that R2 misses. & R2 as proxy (A12); S2--S3 difference; wider range at S3/S4. \emph{Unresolved:} storms missed by R2.\\
Changing travel times & Peak timing error: early (rising river) or late (overbank flow, A10); worst in rare large floods. & Flow-dependent lag attention; high-water weighting. \emph{Unresolved:} floods outside the record.\\
Out-of-record floods & Peaks under-predicted; neural networks extrapolate poorly. & Flag out-of-range inputs; fall back to persistence plus upstream rise.\\
Future and unobserved rain & Rain after $t_0$, or missed by R1/R2, raises S2/S3 unexpectedly. & Inherent at a one-hour horizon; wider range; rain nowcasts requested.\\
Data gaps and batching & Up to 1 h missing before each batch (A4); spikes or stale values. & Masks; training with simulated delays and masking; quality screen.\\
Release deviations & S1, then S2--S4, wrong when operations depart from the plan. & Recorded actuals override the plan; deviation alert (Section 5).\\
Peak smoothing & MSE predicts the conditional mean, which flattens peaks under uncertainty. & High-water weighting; optional quantile output.\\
Short or ageing record & Few floods, so a deep model may not beat simple methods [4]; channel or rule changes age the data. & Small shared-weight model; deploy only if it beats persistence on validation floods; monitor and retrain.\\
\bottomrule
\end{tabularx}
\end{table}

\textbf{Verification of learned delays.} A path does not guarantee that the model uses it; in the companion notebook only one attention head learned to read the source. Attention maps are therefore inspected after training: S3 should attend to S2 at a lag near the travel time, and the lag should shorten as the river rises. If not, the path exists but is unused, and the model is retrained or the design revised.

\begin{table}[H]\centering\small
\caption{Additional information requested from the authority.}
\label{tab:q1requests}
\begin{tabularx}{\linewidth}{@{}>{\raggedright\arraybackslash}p{3.4cm}LL@{}}
\toprule
Request & How it helps & Usable without it?\\
\midrule
Gauge on the unmonitored tributary (temporary is enough) & Removes the main blind spot; calibrates the R2 proxy. & Yes; wider range at S3/S4.\\
Archived release schedules with announcement times & Training on plans as known at $t_0$ (A6). & Partly; actuals as a stand-in overstate plan reliability.\\
Flood alert levels and flood reports & Loss weights; event-based evaluation. & Yes; percentiles instead.\\
Rain radar or nowcasts & Rain expected in the next hour. & Yes; current design.\\
Rating curves or discharge & Water-balance checks; estimate of unmonitored inflow. & Yes.\\
Real-time telemetry and delay logs & Removes pre-batch gaps; realistic delay simulation (A4). & Yes; less sharp before each batch.\\
Channel and floodplain geometry & Identifies overbank flow and longer travel times (A10). & Yes; delay still learned.\\
Catchment boundaries for R1, R2 & Confirms the rain-to-row mapping (A11). & Yes; A11 as stated.\\
Reservoir level and spillway rules & Anticipates unscheduled spills. & Yes; spills seen only at S1.\\
\bottomrule
\end{tabularx}
\end{table}
